% ============================================================
% Autor: Eloi Egea Rada
% Tema: Conclusions, outcomes evaluation, and future work
% ============================================================

\chapter{Conclusions}
\label{ch:conclusions}

This chapter reflects on the lessons learned during development, summarises the
project achievements, and outlines directions for future work. The overall goal ---
to design and implement a reusable, curriculum-aligned penetration testing lab
package for the \textit{Introduction to Cybersecurity} course at TecnoCampus ---
has been achieved. The scope and design evolved iteratively throughout development
in response to infrastructure findings, pilot feedback, and technical constraints,
but all four planned labs were delivered and all stated objectives were substantially
met. The results of the pilot validation and the formal objectives evaluation are
presented in Chapter~\ref{ch:results}.

\section{Summary of Achievements}

The project delivered four progressive lab scenarios built on EVE-NG~\cite{eve_ng}
and deployed on a self-hosted Proxmox~\cite{proxmox} infrastructure:

\begin{itemize}
    \item \textbf{Lab~1 --- Reconnaissance and Enumeration}: A layered network
          topology featuring pfSense~\cite{pfsense}, VyOS~\cite{vyos}, and an
          Ubuntu Server~\cite{ubuntu} target. Students perform passive and active
          reconnaissance using Nmap, DNS enumeration, and service fingerprinting,
          mirroring the PTES~\cite{ptes} initial access phase.

    \item \textbf{Lab~2 --- Web Application Vulnerabilities (SYNAPSE)}: A custom
          two-tier web application aligned with OWASP Top~10~\cite{owasp2021}.
          The exploitation chain --- stored XSS, SQL injection, and YAML deserialization
          --- guides students through a realistic attack scenario that requires
          lateral thinking and multi-step reasoning.

    \item \textbf{Lab~3 --- Incident Response and Log Forensics (HELIX Systems)}:
          An inverted-role lab where students act as security analysts following NIST
          SP~800-61~\cite{nist_sp80061}. Pre-staged log evidence and a reconstructed
          attack timeline develop forensic analysis skills alongside offensive concepts.

    \item \textbf{Lab~4 --- Cryptography and Steganography CTF (CipherStrike)}:
          A three-stage CTF set in the context of a corporate breach at NovaCorp.
          Students recover hidden data from five cryptography challenges on ServerA,
          five steganography challenges on ServerB, and four advanced combined
          challenges on ServerC, whose access is gated on solving key challenges
          from the first two modules. Full credentials and flag inventory are in
          Appendix~J.
\end{itemize}

Each lab is supported by a student guide, a teacher resolution guide with an assessment
rubric, and automation scripts for deployment and environment reset. All documentation
is published in three parallel layers: the formal LaTeX report (this volume) and
Volume~III (Web Documentation, Appendices~A--J).

\section{Lessons Learned}

\paragraph{Infrastructure stability precedes content.}
The migration from VMware to Proxmox early in Phase~0 was the single most impactful
decision of the project. Developing lab content on an unstable hypervisor
would have caused compounding delays. The lesson is general: validate the
execution environment thoroughly before investing in content.

\paragraph{Image management is a project in itself.}
Selecting, importing, and maintaining QEMU images for EVE-NG required significantly
more effort than expected. Naming conventions, permission scripts, and QCOW2
format validation each needed dedicated attention. The \texttt{Iso\_uploader.py}
script emerged from this experience and reduced future onboarding time substantially.

\paragraph{Flag-based progression is effective for self-paced learning.}
The flag model --- where students collect verifiable tokens at each milestone ---
provides immediate feedback and natural progress checkpoints without requiring
teacher presence. The pilot study confirmed that students find the
model motivating and that flags help them identify when they are on the wrong path.

\paragraph{Documentation as a first-class deliverable.}
Producing documentation in three layers (LaTeX, web, src/) added overhead
but proved valuable: the web documentation became the primary reference during
pilot sessions, while the LaTeX volume serves the formal academic record.
Future projects of this type should plan the documentation architecture before
starting implementation.

\paragraph{Time estimates for realistic EVE-NG scenarios are consistently optimistic.}
Each lab required approximately 20--30\% more hours than initially estimated,
mainly due to unexpected node configuration persistence issues and the need to
validate flag reachability after each topology change. Any similar project plan
should include an explicit buffer for integration testing.

\section{Limitations and Future Work}

\subsection{Current Limitations}

\begin{itemize}
    \item \textbf{No LMS integration}: The package is self-contained but not
          integrated with Moodle or any institutional learning management system.
          Assessment rubrics are PDF-based; submission and grading are manual.

    \item \textbf{Single-instance concurrency}: EVE-NG Community Edition runs as a
          single shared instance. If two students or groups need to run labs simultaneously,
          a separate EVE-NG instance is required for each. Scaling to a full classroom
          therefore requires either multiple EVE-NG deployments or a cloud-based setup.
\end{itemize}

\subsection{Directions for Future Work}

\begin{itemize}
    \item \textbf{Expand OWASP coverage}: Lab~2 currently covers five Top~10
          categories. Additional modules could address server-side request forgery
          (SSRF), insecure deserialization, and API security misconfigurations.

    \item \textbf{Moodle integration}: Export assessment rubrics as Moodle-compatible
          grading forms; implement automated flag submission via a simple REST endpoint.

    \item \textbf{Multi-group support}: Evaluate cloud EVE-NG deployment (AWS or
          on-premises scale-out) for concurrent execution by multiple student groups
          within the same course edition.

    \item \textbf{Lab~5 --- Active Directory}: A Windows domain penetration testing
          scenario (Kerberoasting, Pass-the-Hash, lateral movement) to complement
          the existing Linux-centric labs and address OSCP-aligned skill gaps~\cite{oscp_exam}.

    \item \textbf{Automation layer generalisation}: While utility scripts were
          developed for core lab management tasks (ISO upload, bridge configuration,
          connectivity checks), full cross-environment generalisation and robust
          validation pipelines were not achieved within the scope of this project.
          This represents a concrete avenue for future development: extending the
          scripts with environment detection, idempotency checks, and structured
          error reporting would significantly reduce the setup burden for institutions
          adopting the package in heterogeneous infrastructure contexts.
\end{itemize}

\section{Professional Reflection}

The development of this project required the integration of skills acquired
across the Computer Engineering degree --- network administration, operating
system internals, web security, scripting, and technical writing --- within a
single coherent artefact. Designing realistic attack scenarios, implementing
them in a reproducible environment, and documenting them rigorously as teaching
material proved more challenging and more rewarding than initially anticipated.

Working within the constraints of academic infrastructure and a fixed timeline
demanded constant prioritisation. The most durable decisions were those made at the
architectural level early in Phase~0: choosing Proxmox over VMware,
EVE-NG Community Edition over proprietary alternatives, and an open-source-only tool
stack. These choices ensure that the package remains deployable and maintainable
by any institution with comparable hardware, independent of licensing or vendor
relationships.

The project reinforces the view that practical, scenario-based learning is
essential for cybersecurity education, and that well-designed lab environments can
significantly narrow the gap between theoretical knowledge and professional practice.
Delivering a package intended for use in a real course represents the most
meaningful validation of the work.
